\section{Abelian kernels and tight markings}
% BEGIN F1.Abelian
\begin{theorem}[The abelian classification]\label{F1.Abelian}
Under the hypotheses of Theorem~\ref{F1.Structure}, the kernel $P$
is abelian if and only if the marking is tight. Up to duality, all
such markings are the following, one for each $1\leq a\leq b$,
$a+b=m$:
\begin{gather*}
 P=\gen{x}\times\gen{y},\qquad |x|=p^b,\quad |y|=p^a,\\
 x^{r_0}=x,\quad y^{r_0}=y^{-1},\qquad
 x^{r_2}=x^{-1},\quad y^{r_2}=y^{-1},\qquad r_1=xy r_2.
\end{gather*}
The type is $\{2p^a,p^b\}$. For this type the marked group is unique.
\end{theorem}
\begin{proof}
First suppose $P$ is abelian. Write $r_1=ut$ with $u\in P$ and
$1\ne t\in K$. The two-character lemma of~\cite{UF} shows that
$P$ is the direct product of the two nonzero cyclic character
components on which $t$ acts as $-1$. Choose their projected
generators $x,y$ so that $u=xy$.

If $t=r_0r_2$, the two characters on $(r_0,r_2)$ are $(+,-)$ and
$(-,+)$. A direct product calculation gives adjacent orders
$2|y|$ and $2|x|$. The rank-three product bound of~\cite{UF} would
give $8|x||y|\leq |G|=4|x||y|$, which is impossible.
Thus $t$ is an endpoint; by duality take $t=r_2$. The two characters
are now $(+,-)$ and $(-,-)$, exactly as in the displayed model.
For $|x|=p^b$ and $|y|=p^a$ the adjacent orders are
\[
 |r_0r_1|=2p^a,\qquad |r_1r_2|=p^{\max(a,b)}.
\]
The same bound implies $a+\max(a,b)\leq a+b$, hence $a\leq b$.
It is then an equality; the marking is tight.

Conversely, construct the displayed semidirect product for
$1\leq a\leq b$. The two specified automorphisms generate a
Klein four complement, and $r_1$ is an involution because $r_2$
inverts $xy$. The three marked elements are distinct, the endpoints
commute, and the adjacent orders are $2p^a,p^b$.
The elements $u=xy$ and $u^{r_0}=xy^{-1}$ yield $x^2,y^2$;
squaring is invertible on both cyclic odd-order factors. Hence the
marking generates $G$.

Put $H=\gen{r_0,r_1}$ and $J=\gen{r_1,r_2}$.
The image of $H$ in $K$ has order four. Its rotation square
generates $\gen{y}$, and $|H|=4p^a$, so $H\cap P=\gen{y}$.
Similarly $J\cap P=\gen{xy}$ and its image in $K$ is
$\gen{r_2}$. These two odd subgroups meet trivially: if $x^j=1$
then $p^b\mid j$, which also gives $y^j=1$ since $a\leq b$.
The map $H\cap J\longrightarrow K$ is consequently injective,
with image in $\gen{r_2}$. It contains $r_1$, so
$H\cap J=\gen{r_1}$. The rank-three test proves the full
intersection property. The cyclic factor orders, their projected
generators and the two characters determine the marked group,
giving uniqueness among abelian-kernel models.

It remains to show that tightness implies an abelian kernel. By
Theorem~\ref{F1.Structure}, each type parameter is a power of $p$
or twice a power of $p$. The equality $|G|=2k_1k_2=4p^m$
forces exactly one parameter to be even, and the exponents sum to
$m$. We now retain the precise odd/even clause of
Cunningham--Pellicer's classification~\cite[Theorem 1.1(b),(c)]{CP}:
apart from its exceptional even parameter four, a tight regular
polyhedron with one odd and one even type parameter exists only
when the even parameter divides twice the odd parameter, and the
marked polyhedron is unique for such a type. Here the even parameter
is $2p^a$, never four. Thus $a\leq b$, and uniqueness identifies
the marking with the abelian model just constructed. The use of that
theorem includes the retained regular-polyhedron/string-C-group
correspondence. This completes the proof.
\end{proof}
% END F1.Abelian

% BEGIN F1.Correction
\paragraph{The literal formula in the primary edition.}
The display in~\cite[Theorem 1.2, physical p.\ 2]{HFLQ} gives
$x^{r_0}=x$ and $y^{r_0}=y^{-1}$ but writes $r_1=xy r_0$.
Since $x,y$ commute, this element has square $x^2\ne1$.
It is not an involution for the asserted positive odd-prime exponent.
The body of the same edition, Section 4 on physical page 6, uses
$xy r_2$ in the corresponding character case. Theorem~\ref{F1.Abelian}
uses that formula and assigns the larger exponent to $x$.
This is a correction to the specified preprint display, not a recovery
of its impossible literal formula or an assertion about a later
journal edition.
% END F1.Correction
